\pagestyle{plain}
\setcounter{page}{9} % Ajusta a numeração para começar na página 9 conforme o sumário do template

\section{INTRODUÇÃO}
É a apresentação do tema e do problema. Deve delimitar o assunto tratado e contextualizá-lo para esclarecer a existência de uma lacuna a ser estudada (problematização). Conforme \textcite[p. 5]{reis_frota} o estudo poderá necessitar de hipótese(s) que são as possibilidades disponíveis que corroboram com as expectativas da solução do problema de pesquisa. Ela(s) podem ser confirmadas ou refutadas no decorrer do estudo.

O texto deve ser digitado usando a fonte Arial ou Times New Roman, espaçamento entrelinhas 1,5 cm, tamanho 12, inclusive para as informações constantes na capa. Para as citações com mais de três linhas, notas de rodapé, paginação, legendas e fontes das ilustrações e das tabelas deve-se usar a fonte no tamanho 10.

O alinhamento do texto deve ser justificado e as referências no final do trabalho devem ser alinhadas à esquerda, com espaçamento simples (1,0) e um espaço simples entre elas, conforme \textcite{abnt6023}. No que tange a contagem das páginas, devem ser feitas a partir da folha de rosto, mas somente numerada a partir da introdução.

Vale salientar que todos os autores citados no texto devem ter a referência incluída em lista no final no trabalho \cite{abnt14724}.
